\documentclass[11pt]{article}
\usepackage[margin=0.78in]{geometry}
\usepackage{amsmath,amssymb,amsthm,graphicx,microtype}
\usepackage[colorlinks=true,urlcolor=blue,citecolor=blue]{hyperref}

\newtheorem{theorem}{Theorem}
\newtheorem{lemma}{Lemma}
\newcommand{\area}{\operatorname{area}}
\newcommand{\perim}{\operatorname{per}}
\newcommand{\conv}{\operatorname{conv}}

\title{An elementary polygonal proof of the outer-anticircle Bonnesen inequality}
\author{Candidate solution for \texttt{fa\_banach\_001}\\
\texttt{agent\_lane\_03}, model GPT5.6}
\date{26 August 2026}

\begin{document}
\maketitle

\begin{abstract}
Martini and Swanepoel asked for an elementary polygonal proof of their
inequality
\[
 \area(C)+\area(\sigma I)\leq \sigma\perim(C),
\]
where $\sigma I$ is a smallest anticircle containing a planar convex body
$C$.  We give such a proof.  On a common polygonal normal fan, the difference
between the two sides is the quadratic area form applied to the support
deficit $h_I-h_C$.  Minimum circumscription supplies zero deficits whose
angular gaps are at most $\pi$.  An exact two-angle completion-of-squares
identity makes the quadratic form nonpositive on every such gap.  A common
circumscribed-polygon approximation finishes the proof for arbitrary convex
bodies.  No mixed-area inequality is used.
\end{abstract}

\section{Source request}

Martini and Swanepoel, \emph{Antinorms and Radon curves},
arXiv:math/0409200 (published in \emph{Aequationes Mathematicae} 72 (2006),
110--138), prove in Theorem~15 that, when $\sigma I$ is a smallest
anticircle containing $C$,
\[
 \area(C)+\area(\sigma I)\leq\sigma\perim(C). \tag{13}
\]
At the top of PDF page 19 they ask for an elementary proof involving
polygons, rather than the proof through Blaschke's inequality \cite{MS}.

\begin{center}
 \includegraphics[width=0.94\linewidth]{figures/open_problem_crop.png}
\end{center}

\section{Idea of the Proof}

Minimum circumscription forces the common supporting directions of $C$ and
$I$ to surround the origin; two or three of them therefore cut the full
normal circle into arcs of length at most $\pi$.  For polygons on a common
normal fan, area is an explicit quadratic form $Q$ in the support numbers,
and perimeter is its polarization against the support vector of $I$.  Thus
the desired deficit is simply $Q(h_I-h_C)$.  The support deficit vanishes at
the selected contacts.  On each intervening arc, $-2Q$ is a discrete
Dirichlet energy, and a two-angle completion of squares makes that energy
nonnegative.  Circumscribed polygonal approximation retains the balanced
contacts and passes the result to arbitrary convex bodies.

\section{Statement and normalization}

Fix linear coordinates in which the given oriented area form is the standard
determinant, and use the associated Euclidean scalar product only to
parametrize supporting directions.  Let $I$ be the unit anticircle of a
Minkowski norm.  We use the elementary Steiner identity
\begin{equation}\label{eq:steiner}
 \area(A+tI)=\area(A)+t\perim(A)+t^2\area(I),\qquad t\geq0.
\end{equation}
For polygons this follows directly by adding the edge strips and corner
triangles; equivalently, $\perim(A)$ is twice the polarization of the area
quadratic form evaluated at $A$ and $I$.  This identity is also established
in the antinorm framework of the source paper \cite{MS}.

Translate the smallest containing anticircle to have center $0$ and dilate
by $1/\sigma$.  It is enough to prove the following normalized statement.

\begin{theorem}\label{thm:main}
Let $C\subseteq I$, and assume that no translate of $rI$ contains $C$ for
any $r<1$.  Then
\[
 \area(C)+\area(I)\leq\perim(C).
\]
\end{theorem}

We first isolate the only consequence of minimum circumscription that will
be needed.  For a convex body $A$, write $h_A(u)=\max_{x\in A}\langle
x,u\rangle$.

\begin{lemma}[balanced contact normals]\label{lem:contact}
Under the hypotheses of Theorem~\ref{thm:main}, there are two or three unit
directions $u_j$ such that
\[
 h_C(u_j)=h_I(u_j)
 \quad\text{and}\quad
 0\in\conv\left\{\frac{u_j}{h_I(u_j)}\right\}.
\]
In particular, in their cyclic order, the angular gap between consecutive
selected directions is at most $\pi$.
\end{lemma}

\begin{proof}
For a proposed center $x$, the least radius of a translate centered at $x$
which contains $C$ is
\[
 R(x)=\max_{u\in S^1}
 \frac{h_C(u)-\langle x,u\rangle}{h_I(u)}.
\]
The hypothesis says that $R$ has its minimum value $1$ at $x=0$.  Its active
directions there are exactly
\[
 \mathcal A=\{u:h_C(u)=h_I(u)\}.
\]
If $0$ were outside
$\conv\{u/h_I(u):u\in\mathcal A\}$, strict separation and compactness would
give a vector $v$ such that $\langle v,u\rangle/h_I(u)>0$ uniformly on the
active set.  By continuity, the same strict decrease of the quotient under
$x=\varepsilon v$ holds on a neighborhood $U$ of the active set.  On the
compact complement $S^1\setminus U$, the quotient at $x=0$ is at most
$1-\eta$ for some $\eta>0$, so it remains below $1$ for sufficiently small
$\varepsilon>0$.  Consequently $R(\varepsilon v)<1$, contradicting
minimality.  Carath\'eodory's theorem in the plane now selects at most three
active directions.  A set of directions whose convex hull contains $0$
cannot lie in an open semicircle, which is the final assertion.
\end{proof}

\section{The discrete inequality}

For $0<\alpha<\pi$, define
\begin{equation}\label{eq:ealpha}
 E_\alpha(x,y)=
 \frac{\cos\alpha(x^2+y^2)-2xy}{\sin\alpha}.
\end{equation}

\begin{lemma}[discrete Dirichlet inequality]\label{lem:energy}
Let $\alpha_0,\ldots,\alpha_{m-1}>0$ and
$\Theta=\sum_i\alpha_i\leq\pi$.  If $x_0=x_m=0$, then
\[
 \sum_{i=0}^{m-1}E_{\alpha_i}(x_i,x_{i+1})\geq0.
\]
\end{lemma}

\begin{proof}
For $a,b>0$ with $a+b<\pi$, direct expansion and
$\cot a+\cot b=\sin(a+b)/(\sin a\sin b)$ give the exact identity
\begin{align}\label{eq:merge}
 &E_a(x,z)+E_b(z,y)-E_{a+b}(x,y)\\
 &\quad=
 \frac{\sin(a+b)}{\sin a\sin b}
 \left(z-\frac{x\sin b+y\sin a}{\sin(a+b)}\right)^2\geq0.\nonumber
\end{align}
If $\Theta<\pi$, repeatedly merge adjacent intervals using
\eqref{eq:merge}.  The sum is bounded below by
$E_\Theta(x_0,x_m)=E_\Theta(0,0)=0$.  If $\Theta=\pi$, replace every
$\alpha_i$ by $t\alpha_i$, apply the strict case for $0<t<1$, and let
$t\uparrow1$.
\end{proof}

\section{Polygonal proof}

Suppose first that $C$ and $I$ are polygons.  Take a common cyclic normal
fan containing every side normal of both polygons, the contact normals from
Lemma~\ref{lem:contact}, and enough extra directions that adjacent angles
are below $\pi$.  Write
\[
 u_i=(\cos\theta_i,\sin\theta_i),\qquad
 \alpha_i=\theta_{i+1}-\theta_i\in(0,\pi),
\]
with indices read cyclically, and let $h_i$ be the support number of a
polygon at $u_i$.  Intersecting the supporting half-planes recovers the
polygon; an inserted direction merely gives a possibly degenerate side.

The intersection of the two adjacent support lines gives the standard
support-number area formula
\begin{equation}\label{eq:Q}
 Q(h)=
 \sum_i\frac{h_i h_{i+1}}{\sin\alpha_i}
 -\frac12\sum_i h_i^2
 \bigl(\cot\alpha_{i-1}+\cot\alpha_i\bigr).
\end{equation}
Indeed, substituting those line intersections into the shoelace formula
gives \eqref{eq:Q} term by term.  Thus $Q(h_A)=\area(A)$.  Let $B$ be the
polarization of $Q$.  Support numbers add under Minkowski addition, so
\eqref{eq:steiner} (or direct strip decomposition) gives
\begin{equation}\label{eq:perpol}
 \perim(C)=2B(h_C,h_I).
\end{equation}

Set $c_i=h_C(u_i)$, $k_i=h_I(u_i)$ and $d_i=k_i-c_i\geq0$.  Quadraticity and
\eqref{eq:perpol} show that the desired inequality is exactly
\begin{equation}\label{eq:deficitQ}
 \area(C)+\area(I)-\perim(C)
 =Q(c)+Q(k)-2B(c,k)=Q(k-c)=Q(d)\leq0.
\end{equation}
Expanding \eqref{eq:Q} in edge form gives
\begin{equation}\label{eq:energyQ}
 -2Q(d)=\sum_i E_{\alpha_i}(d_i,d_{i+1}).
\end{equation}

At each selected contact normal, $d_i=0$.  Those zero entries split the
cyclic sum in \eqref{eq:energyQ} into paths.  By
Lemma~\ref{lem:contact}, the total angle of each path is at most $\pi$.
Lemma~\ref{lem:energy} makes every path energy nonnegative.  Hence
$-2Q(d)\geq0$, proving \eqref{eq:deficitQ} and Theorem~\ref{thm:main} for
polygons.

\section{Passage to arbitrary convex bodies}

Choose finite symmetric direction sets $\mathcal U_n\subset S^1$ whose mesh
tends to zero and which all contain the selected contact directions from
Lemma~\ref{lem:contact}.  Define simultaneous circumscribed polygons
\[
 C_n=\bigcap_{u\in\mathcal U_n}
 \{x:\langle x,u\rangle\leq h_C(u)\},\qquad
 I_n=\bigcap_{u\in\mathcal U_n}
 \{x:\langle x,u\rangle\leq h_I(u)\}.
\]
Then $C_n\subseteq I_n$, and both sequences converge in Hausdorff distance
to $C$ and $I$.  The chosen contact constraints remain equalities for both
polygons.  Moreover $I_n$ is still a minimum-radius containing homothet for
$C_n$: if $C_n\subseteq x+rI_n$, take convex coefficients $\lambda_j$ from
Lemma~\ref{lem:contact} and apply the support inequalities at the selected
contacts to obtain
\[
 1\leq r+\left\langle x,
 \sum_j\lambda_j\frac{u_j}{h_I(u_j)}\right\rangle=r.
\]
The polygonal argument above is a purely convex statement about a pair
$(C,I)$ with $I$ centrally symmetric: its right side is $2B(C,I)$, and the
only hypotheses used are containment and minimum homothety.  Since
$\mathcal U_n$ is symmetric, $I_n$ is centrally symmetric, so that argument
applied to $(C_n,I_n)$ gives
\[
 \area(C_n)+\area(I_n)\leq 2B(C_n,I_n).
\]
Area and its polarization are continuous under planar Hausdorff convergence.
Letting $n\to\infty$ and using \eqref{eq:steiner} yields
\[
 \area(C)+\area(I)\leq2B(C,I)=\perim(C).
\]
This proves the normalized theorem.  Undoing the translation and the
dilation by $1/\sigma$ gives
\[
 \boxed{\area(C)+\area(\sigma I)\leq\sigma\perim(C)}.
\]

\section*{Novelty and scope}

The theorem is classical; only the elementary proof was requested.  The
four run indexes, exact-title citations in the local parsed arXiv corpus, and
current exact-phrase searches for the missing polygonal proof were checked on
26 August 2026.  They located mixed-area proofs and papers citing
arXiv:math/0409200 for other antinorm facts, but no later proof matching the
argument above.  This is bounded novelty evidence, not a guarantee of
publication novelty.

\section*{Verification and review recommendation}

The accompanying script \texttt{code/verify\_discrete\_energy.py} checks the
merge identity, random zero-endpoint path energies, the identity
$Q=-E/2$, and the support-number area formula against shoelace area.  With
20,000 trials and seed 20260826 all checks pass.  These computations are
consistency checks, not part of the proof.  Human review should concentrate
on the support-number formula when redundant normals are inserted, the
balanced-contact argument for translated homothets, and the simultaneous
circumscribed-polygon limit.

\begin{thebibliography}{9}
\bibitem{MS}
H.~Martini and K.~J.~Swanepoel,
\emph{Antinorms and Radon curves},
Aequationes Math. \textbf{72} (2006), 110--138;
arXiv:math/0409200,
\url{https://arxiv.org/abs/math/0409200}.
\end{thebibliography}

\end{document}
